% ══════════════════════════════════════════════════════════════════
% Byzantine-Based Blockchain Consensus for EV-to-Grid Energy Trading:
% A Static-Temporal Security Evaluation Framework
% ══════════════════════════════════════════════════════════════════
\documentclass[10pt,twocolumn,journal]{IEEEtran}

% ─── Packages ────────────────────────────────────────────────────
\usepackage{amsmath,amssymb,amsthm}
\usepackage{booktabs,multirow,array,colortbl}
\usepackage{graphicx,subcaption}
\graphicspath{{figures/}}
\usepackage{geometry,xcolor,enumitem}
\definecolor{ieee}{RGB}{0,83,159}
\usepackage[colorlinks=true, allcolors=ieee]{hyperref}
\usepackage{cite}
\usepackage{url}

\geometry{margin=0.75in}
\newtheorem{theorem}{Theorem}
\newtheorem{definition}{Definition}
\newtheorem{proposition}{Proposition}

% ═════════════════════════════════════════════════════════════════
\title{\textbf{Byzantine-Based Blockchain Consensus for EV-to-Grid Energy Trading:\\A Static-Temporal Security Evaluation Framework}\\[8pt]
\large Evaluating Cyber-Physical Security Improvements in Advanced Metering Infrastructure}

\author{\textbf{Atharv~Manojkumar~Shukla}\\
\vspace{0.15cm}
\small \emph{Co-Authors / Project Guides:} \textbf{Prof.~Uday~Suryavanshi} and \textbf{Dr.~Sunny~Kumar}%
\thanks{A. M. Shukla is a 3rd year B.Tech Undergraduate Student with the Department of Computer Science and Engineering, Indian Institute of Information Technology (IIIT), Nagpur 441108, India (e-mail: atharv@iiitn.ac.in).}%
\thanks{Prof. Uday Suryavanshi and Dr. Sunny Kumar are Co-Authors \& Project Guides with the E-MC$^2$ Laboratory, Department of Electrical Engineering, Veermata Jijabai Technological Institute (VJTI), Mumbai 400019, India (e-mail: usuryavanshi@ee.vjti.ac.in, skumar@ee.vjti.ac.in).}}

\markboth{Research Project Report --- E-MC$^2$ Laboratory, VJTI Mumbai}%
{Shukla \MakeLowercase{\textit{et al.}}: Byzantine-Based Blockchain Consensus for EV-to-Grid Energy Trading}

\begin{document}
\maketitle

% ═══════════════════════════════════════════════════════════════
% ABSTRACT
% ═══════════════════════════════════════════════════════════════
\begin{abstract}
This paper quantifies the cybersecurity improvement that Byzantine Fault Tolerant (BFT) blockchain consensus mechanisms introduce to Electric Vehicle (EV) energy trading networks under the \emph{Static–Temporal Security Framework (STSF)}. Moving beyond simple protocol benchmarking, we focus on the central question: \emph{How much does blockchain improve the security of an EV Energy Trading System compared to the same system without blockchain?} Using the IEEE 33-bus distribution system with 50 EVs and 3,854 physical sensors as our reference, we evaluate security improvements using Sheikh et al.'s foundational four-component probabilistic attack model, demonstrating that under baseline model assumptions, blockchain reduces total attack probability from $P_{\text{TA}} \approx 0.005$ to $P_{\text{TAb}} \approx 10^{-173}$---a static security gain of approximately 170 orders of magnitude under the Sheikh formulation. We complement this with a probabilistic evaluation across 12 distinct cyber-physical attack vectors (including Sybil, Byzantine node, and DDoS attacks) to illustrate the breadth of defensive coverage. We then prove that consensus validation latency creates a temporal vulnerability window ($P_{\text{temporal}}$) governed by a Poisson arrival process. Pipelined sub-second protocols (e.g., Tower BFT, RVR) preserve over 92\% of static security gains ($P_{\text{secure}} \ge 0.927$), whereas classical consensus protocols (e.g., $OM(m)$) lose up to 99\% of security benefits due to latency exposure ($P_{\text{secure}} = 0.012$). Finally, we corroborate our analytical formulations across $10^6$ Monte Carlo simulation trials to within 0.5\% error with 95\% confidence intervals.
\end{abstract}

% ═══════════════════════════════════════════════════════════════
\section{Introduction}
\label{sec:introduction}
% ═══════════════════════════════════════════════════════════════

The modernization of electrical grids into cyber-physical smart distribution networks has enabled bidirectional energy and data flows. In particular, peer-to-peer (P2P) energy trading between Electric Vehicles (EVs) and Advanced Metering Infrastructure (AMI) provides grid balancing and ancillary services. However, this decentralized structure exposes the system to cyber-physical threats, including False Data Injection (FDI), Man-in-the-Middle (MitM), and Denial-of-Service (DoS) attacks.

While blockchain consensus has been proposed to secure these networks, existing literature often treats security as an instantaneous, static property or focuses solely on consensus performance benchmarking. This paper addresses a more fundamental research question:

\begin{quote}
\textbf{``How much does blockchain improve the security of an EV Energy Trading System compared to the same system without blockchain, and how does the choice of consensus algorithm affect this difference?''}
\end{quote}

To address this, we structure the evaluation into two parts. In the first part, we model the system without blockchain, derive the success probabilities for 12 core cyber-attacks, and compare them directly against a blockchain-protected baseline. In the second part, we evaluate how the temporal exposure window of different consensus algorithms affects the realized security gains, using throughput, latency, communication cost, and scalability.

Specifically, relative to the baseline study by Sheikh et al.~\cite{sheikh2020}, the core contributions proposed and evaluated in this paper are:
\begin{itemize}[leftmargin=*]
    \item \textbf{Localized Sensor Subset Formulation:} Rather than utilizing global network sensor counts (which lead to analytical underflows), we propose a localized critical-subset formulation ($m=10$) to model cyber-physical attacks realistic to smart grid operations.
    \item \textbf{Novel Comparative Static-Temporal Security Model (STSF):} We propose a first-order probabilistic security model that couples static cryptosystem compromise bounds with dynamic consensus-dependent exposure latency ($P_{\text{secure}}$).
    \item \textbf{Rigorous With-vs-Without Blockchain Evaluation:} We map and formulate the 12 cyber-physical attack success models explicitly under both centralized coordinator architectures and permissioned BFT blockchain environments.
    \item \textbf{Consensus Selection Roadmap:} We supply an operational decision roadmap by compiling theoretical message complexities and normalizing heterogeneous performance metrics (throughput, latency, and communication overhead).
\end{itemize}

% ═══════════════════════════════════════════════════════════════
% PART I — APPLICATION SECURITY ANALYSIS (PRIMARY CONTRIBUTION)
% ═══════════════════════════════════════════════════════════════

\section*{Part I: Application Security}
\section{Without Blockchain}
\label{sec:part1}

\emph{This section addresses the paper's primary research objective: quantifying how much blockchain improves the security of the EV Energy Trading application. The comparison is strictly between the centralized (no-blockchain) and decentralized (blockchain-protected) configurations.}

\subsection{System Model}

We model an EV-integrated IEEE 33-bus distribution system (33 buses, 32 branches, peak load of 3,715\,kW) with 50 EVs serving as prosumers. This yields $n=51$ participant validator nodes. The system contains $n_{\text{sen}} = 3,854$ physical sensors:
\begin{equation}
n_{\text{sen}} = n_{\text{sen}}^{\text{DN}} (129) + n_{\text{sen}}^{\text{EV}} (3,725) = 3,854
\end{equation}

\subsection{Sensor Count Dominance Critique}

Prior models of smart meter security express sensor compromise probability ($P_{\text{SA}}$) and communication link compromise probability ($P_{\text{CA}}$) as:
\begin{equation}
P_{\text{SA}} = P_{\text{CA}} = x^{n_{\text{sen}}}
\end{equation}
where $x \in [0.9, 0.999]$ is the per-sensor security level. At $x = 0.95$ and $n_{\text{sen}}=3,854$, $x^{3854} \approx 10^{-86}$. This causes the sensor and communication terms to underflow to zero, making system security appear entirely independent of sensor vulnerability.

In reality, an attacker only needs to compromise a localized subset of critical sensors ($m$) at key bus nodes (e.g., feeder heads or lateral branch points) to inject false data or disrupt charging operations. We define $m = 10$, restoring model sensitivity:
\begin{equation}
P_{\text{SA}}^{(10)} = P_{\text{CA}}^{(10)} = x^{10} = 0.95^{10} \approx 0.599
\end{equation}

\subsection{Justification and Causal Analysis of the 12 Cyber-Physical Attacks}

Without blockchain, the system relies on centralized coordinators or weak client-server protocols. The baseline success probabilities ($P_{\text{atk}}^j$) for 12 cyber-attacks are justified below:

\begin{enumerate}[leftmargin=*]
    \item \textbf{Sensor Compromise ($P_{\text{SA}}$):} $P_{\text{SA}} = x^m = 0.599$. *Causal mechanism:* Attacker alters local smart meter sensor calibrations without cross-verification.
    \item \textbf{False Data Injection ($P_{\text{FDI}}$):} $P_{\text{FDI}} = x^m = 0.599$. *Causal mechanism:* Altered meter values pass to coordinator state estimation undetected.
    \item \textbf{Communication Compromise ($P_{\text{CA}}$):} $P_{\text{CA}} = x^m = 0.599$. *Causal mechanism:* Intercepting unencrypted transmission channels.
    \item \textbf{Man-in-the-Middle ($P_{\text{MitM}}$):} $P_{\text{MitM}} = y = 0.05$. *Causal mechanism:* TLS session hijacking or DNS spoofing on legacy grid protocols.
    \item \textbf{Replay Attack ($P_{\text{Replay}}$):} $P_{\text{Replay}} = z = 0.15$. *Causal mechanism:* Capturing and replaying valid energy credit transmissions.
    \item \textbf{Sybil Attack ($P_{\text{Sybil}}$):} $P_{\text{Sybil}} = 1.0$. *Causal mechanism:* Open registration allows virtual identity creation.
    \item \textbf{Denial of Service ($P_{\text{DoS}}$):} $P_{\text{DoS}} = p_{\text{dos}} = 0.20$. *Causal mechanism:* Single-node port flooding.
    \item \textbf{Distributed DoS ($P_{\text{DDoS}}$):} $P_{\text{DDoS}} = p_{\text{ddos}} = 0.35$. *Causal mechanism:* Botnet ingress bandwidth exhaustion.
    \item \textbf{Byzantine Validator Attack ($P_{\text{Byz}}$):} $P_{\text{Byz}} = 1.0$. *Causal mechanism:* Absolute central coordinator Single Point of Failure (SPOF).
    \item \textbf{Key Compromise ($P_{\text{Key}}$):} $P_{\text{Key}} = p_{\text{key}} = 0.01$. *Causal mechanism:* Private key leakage from central server filesystem.
    \item \textbf{SCADA Compromise ($P_{\text{SCADA}}$):} $P_{\text{SCADA}} = 0.01$. *Causal mechanism:* Database SQL injection at central SCADA host.
    \item \textbf{Receiver Compromise ($P_{\text{R}}$):} $P_{\text{R}} = 0.01$. *Causal mechanism:* Direct actuator relay firmware overrides.
\end{enumerate}

\subsection{Total Probability of Security Without Blockchain}

Assuming parallel failure modes, the system remains secure only if all attack vectors fail:
\begin{equation}
P_{\text{secure}}^{\text{no\_BC}} = \prod_{j=1}^{12} (1 - P_{\text{atk}}^j) = 0.0
\end{equation}
Since $P_{\text{Sybil}} = 1.0$ and $P_{\text{Byz}} = 1.0$ under centralized architectures, these single points of failure drive overall security to zero. Excluding SPOFs to evaluate non-catastrophic vectors yields $P_{\text{secure, limited}}^{\text{no\_BC}} \approx 0.0263$ (only 2.6\% safe without blockchain).

\section{With Blockchain}
\subsection{Quantitative Security Comparison Using Sheikh's Four-Component Model}
Using Sheikh et al.'s four-component probabilistic model under baseline parameters ($N_{\text{sen}} = 3,854$), applying permissioned BFT consensus reduces attack compromise probability from $P_{\text{TA}} \approx 0.005$ to $P_{\text{TAb}} \approx 4.91 \times 10^{-173}$---a static security gain of 170 orders of magnitude within the analytical model assumptions.

\subsection{12 Cyber-Physical Attack Mitigation Under Blockchain}
\begin{itemize}[leftmargin=*]
    \item \textbf{Sensor \& FDI ($P_{\text{Sensor}} = P_{\text{FDI}} = x^{2m} \approx 0.358$):} Requires both physical sensor override ($x^m$) and private key extraction ($x^m$).
    \item \textbf{Communication Hijack ($P_{\text{Comm}} = x^{2k_1} \approx 0.0099$):} $k_1 = \frac{m(m-1)}{2} = 45$ directional links require dual routing and session key compromise.
    \item \textbf{Man-in-the-Middle ($P_{\text{MitM}} = y^2 = 0.0025$):} Session hijacking combined with cryptographic key extraction.
    \item \textbf{Replay Attack ($P_{\text{Replay}} \to 0.0$):} Completely eliminated via block nonces and cryptographic timestamps.
    \item \textbf{Sybil \& Byzantine ($P_{\text{Sybil}} = P_{\text{Byz}} = \sum_{k=17}^{51} \binom{51}{k} p_c^k (1-p_c)^{51-k} \approx 1.54 \times 10^{-10}$):} Lowered from $1.0$ to $10^{-10}$ via $3f+1$ permissioned BFT quorums.
    \item \textbf{DoS \& DDoS ($P_{\text{DoS}} = p_{\text{dos}} \frac{f+1}{n} \approx 0.067, P_{\text{DDoS}} = p_{\text{ddos}} \frac{f+1}{n} \approx 0.117$):} Multi-node P2P routing requires flooding $f+1$ validators.
    \item \textbf{Key Compromise ($P_{\text{Key}} \approx 9.85 \times 10^{-6}$):} Protected via Shamir $(3,5)$ Threshold Secret Sharing.
    \item \textbf{SCADA \& Receiver ($P_{\text{SCADA}} \approx 5.99 \times 10^{-21}, P_{\text{Receiver}} \approx 5.99 \times 10^{-7}$):} Multi-signature ledger checkpointing filters illegal command execution.
\end{itemize}

% ═══════════════════════════════════════════════════════════════
% PART II — TEMPORAL CONSENSUS EVALUATION & MONTE CARLO
% ═══════════════════════════════════════════════════════════════

\section*{Part II: Temporal Consensus & Monte Carlo Verification}
\section{Complexity-Derived Latency & Poisson Attack Model}

We establish a closed-form derivation expressing consensus validation latency $\tau(n)$ as a function of message complexity $M(n)$, payload size $S$, bandwidth $B$, and node service rate $\mu$:
\begin{equation}
\tau(n) = \tau_{\text{prop}} + \frac{M(n) \cdot S}{B} + \frac{\lambda \cdot S}{\mu - \lambda \cdot S}
\end{equation}
Equation (8) combines physical propagation delay, message complexity transmission delay, and queueing delay, where the queueing delay term is derived from classical $M/M/1$ queueing theory (Kleinrock 1975). We note that Lamport's $OM(m)$ algorithm ($\tau = 43.35\text{s}$) serves solely as the theoretical origin of Byzantine agreement and is included to illustrate historical consensus evolution rather than as a practical deployment candidate.

Under a Poisson attack process ($\lambda = 20\text{ attacks/s}$), time-dependent system survival probability $P_{\text{secure}}(\lambda, \tau)$ is:
\begin{equation}
P_{\text{secure}}(\lambda, \tau) = (1 - P_{\text{TAb}}) \cdot \exp\left(-\lambda \cdot \tau(n) \cdot P_{\text{TA}}\right)
\end{equation}

Sub-second Generation-4 protocols (RVR: $P_{\text{secure}} = 0.9404, \tau = 200\text{ms}$; Tower BFT: $P_{\text{secure}} = 0.9272, \tau = 242.9\text{ms}$) preserve over 92\% of static security gains, whereas Classic PBFT ($\tau = 7.65\text{s}$) suffers severe security decay ($P_{\text{secure}} = 0.4421$).

\section{Monte Carlo Verification & Reproducibility}

Across $N = 10^6$ independent Monte Carlo simulation trials (`seed = 42`), empirical sample means corroborate analytical closed-form equations to within 0.5\% relative error with 95\% confidence intervals:
\begin{itemize}[leftmargin=*]
    \item $P_{\text{SA}}$ Analytical: 0.59871 vs. Monte Carlo: 0.59888 ($\text{Error} = 0.028\%$).
    \item $P_{\text{SA, BFT}}$ Analytical: 0.35853 vs. Monte Carlo: 0.35867 ($\text{Error} = 0.040\%$).
    \item $P_{\text{MitM, BFT}}$ Analytical: 0.00250 vs. Monte Carlo: 0.00251 ($\text{Error} = 0.400\%$).
    \item $P_{\text{Key, BFT}}$ Analytical: $9.851 \times 10^{-6}$ vs. Monte Carlo: $9.848 \times 10^{-6}$ ($\text{Error} = 0.030\%$).
\end{itemize}

\section{Conclusion}
This paper presented the \emph{Static–Temporal Security Framework (STSF)}, quantifying the security improvements of BFT blockchain consensus in EV energy trading networks. Within baseline model assumptions, BFT consensus achieves a 170 order-of-magnitude static security gain, eliminates single-point-of-failure risks ($P_{\text{SPOF}} \to 10^{-10}$), and neutralizes replay vectors. Furthermore, sub-second pipelined consensus protocols (Tower BFT, RVR) preserve over 92\% of static security gains under real-time Poisson attack traffic.

% ═══════════════════════════════════════════════════════════════
% REFERENCES
% ═══════════════════════════════════════════════════════════════
\begin{thebibliography}{99}

\bibitem{sheikh2020}
A.~Sheikh, A.~Ahmadinia, and R.~Kianoush, ``Secured energy trading using Byzantine-based blockchain consensus in smart grid,'' \emph{IEEE Access}, vol.~8, pp.~156\,832--156\,845, 2020.

\bibitem{lamport1982}
L.~Lamport, R.~Shostak, and M.~Pease, ``The Byzantine generals problem,'' \emph{ACM Trans. Program. Lang. Syst.}, vol.~4, no.~3, pp.~382--401, 1982.

\bibitem{castro1999}
M.~Castro and B.~Liskov, ``Practical Byzantine fault tolerance,'' in \emph{Proc. OSDI}, 1999, pp.~173--186.

\bibitem{li2020_blockchain}
Z.~Li, J.~Kang, R.~Yu, D.~Ye, Q.~Deng, and Y.~Zhang, ``Consortium blockchain for secure energy trading in industrial IoT,'' \emph{IEEE Trans. Inf. Inf.}, vol.~14, no.~8, pp.~3690--3700, 2018.

\bibitem{baumeister2010}
T.~Baumeister, ``Literature review on smart grid cyber security,'' Collaborative Softw. Develop. Lab. at Univ. Hawaii, Honolulu, HI, USA, Tech. Rep., 2010.

\bibitem{sridhar2012}
S.~Sridhar, A.~Hahn, and M.~Govindarasu, ``Cyber-physical system security for the electric power grid,'' \emph{Proc. IEEE}, vol.~100, no.~1, pp.~210--224, 2012.

\bibitem{wang2013}
W.~Wang and Z.~Lu, ``Cyber security in the smart grid: Survey and challenges,'' \emph{Computer Networks}, vol.~57, no.~5, pp.~1344--1371, 2013.

\bibitem{yan2012}
Y.~Yan, Y.~Qian, H.~Sharif, and D.~Tipper, ``A survey on cyber security for smart grid communications,'' \emph{IEEE Commun. Surveys Tuts.}, vol.~14, no.~4, pp.~998--1010, 2012.

\bibitem{ding2024}
X.~Ding \emph{et al.}, ``CE-PBFT: An efficient cluster-based PBFT consensus algorithm for microgrid energy trading,'' \emph{IEEE Trans. Smart Grid}, vol.~15, no.~2, pp.~1820--1832, 2024.

\end{thebibliography}

\end{document}
